%=======================================================================
\chapter{Extensions exploratoires : Deep Learning et VLM}

Au-delà du pipeline ML des chapitres précédents, deux notebooks complémentaires ont été
réalisés~: \texttt{Deep\_learning.ipynb} (benchmark CNN et modèles hybrides, Phase~3) et
\texttt{VLM.ipynb} (fine-tuning d'un modèle vision-langage avec LoRA). Cette section en
résume le protocole, les figures clés et les résultats principaux.

% -----------------------------------------------------------------------
\section{Benchmark Deep Learning (\texttt{Deep\_learning.ipynb})}

\subsection{Protocole}

Le notebook implémente une \textbf{Phase~3} structurée en trois niveaux, sur le même jeu de
données (\NTrain{} / \NTest{}) et les mêmes caractéristiques manuelles optimisées (Phase~1)~:

\begin{itemize}[leftmargin=1.4em,itemsep=2pt]
  \item \textbf{Tier~0 --- CNN purs}~: VGG16, ResNet50, EfficientNet-B0 et ViT-B/16
        (transfert ImageNet, tête de classification entraînée).
  \item \textbf{Tier~1 --- Hybrides}~: ResNet18 + une branche manuelle (GLCM, LBP, DWT ou HOG)
        concaténée aux embeddings CNN (\texttt{HybridClassifier}).
  \item \textbf{Tier~2 --- Fusion profonde}~: ResNet18 + les quatre branches manuelles fusionnées
        par concaténation (\texttt{MultiBranchFusion}) ou attention apprise
        (\texttt{AttentionFusion}).
\end{itemize}

Chaque modèle est évalué en \emph{hold-out} (split 80/20 stratifié) et en validation croisée
5-fold. Les métriques sont exportées dans \texttt{results/phase3\_all\_metrics.csv}.

\subsection{Résultats principaux}

Le Tableau~\ref{tab:phase3-dl} et la Figure~\ref{fig:phase3-tier} rapportent les F1-Macro sur
le split de test. Le meilleur modèle pur est \textbf{ResNet50} (F1 $= 0{,}935$), devant
EfficientNet-B0 ($0{,}926$) et VGG16 ($0{,}910$). ViT-B/16 sous-performe sur ce petit corpus
($0{,}842$).

Les hybrides Tier~1 surpassent nettement les branches ML correspondantes de la Phase~2~: par
exemple, \texttt{cnn\_hog} atteint F1 $= 0{,}934$ contre $0{,}903$ pour SVM + D (HOG), et
\texttt{waveletcnn} atteint $0{,}917$ contre $0{,}760$ pour Extra Trees + C (DWT).

En Tier~2, la fusion profonde n'apporte qu'un gain marginal sur Full(opt) de la Phase~2~:
\texttt{multibranch\_fusion} ($0{,}923$) dépasse SVM + Full(opt) ($0{,}920$) de
$+0{,}002$, tandis que \texttt{attention\_fusion} reste en dessous ($0{,}901$).

\begin{table}[htbp]
  \centering
  \caption{Phase~3 --- meilleurs résultats par tier (split de test).}
  \label{tab:phase3-dl}
  \begin{tabular}{@{}clcc@{}}
    \toprule
    \thd{Tier} & \thd{Modèle} & \thd{F1-Macro} & \thd{Accuracy} \\
    \midrule
    0 & ResNet50          & 0,935 & 0,941 \\
    0 & EfficientNet-B0   & 0,926 & 0,933 \\
    0 & VGG16             & 0,910 & 0,915 \\
    1 & CNN + HOG         & 0,934 & 0,940 \\
    1 & LBPNet            & 0,934 & 0,940 \\
    2 & MultiBranchFusion & 0,923 & 0,930 \\
    \midrule
    \multicolumn{2}{@{}l}{\textit{Réf. Phase~2 --- SVM + Full(opt)}} &
      \textbf{0,920} & \textbf{0,920} \\
    \bottomrule
  \end{tabular}
\end{table}

\begin{figure}[htbp]
  \centering
  \includegraphics[width=0.92\textwidth]{phase3_tier_comparison}
  \caption{Comparaison des F1-Macro par modèle et par tier (split de test).}
  \label{fig:phase3-tier}
\end{figure}

La Figure~\ref{fig:phase3-roc} montre les courbes ROC one-vs-rest pour l'ensemble des dix
modèles Phase~3. ResNet50 et les hybrides Tier~1 dominent sur les trois classes, confirmant la
bonne séparabilité des tumeurs dans l'espace des représentations apprises.

\begin{figure}[htbp]
  \centering
  \includegraphics[width=\textwidth]{phase3_roc_all_models}
  \caption{Courbes ROC (one-vs-rest) --- benchmark Phase~3 complet.}
  \label{fig:phase3-roc}
\end{figure}

La matrice de confusion de ResNet50 (Figure~\ref{fig:phase3-cm-resnet50}) confirme un faible
taux d'erreurs inter-classes, avec une légère confusion entre gliome et méningiome.

\begin{figure}[htbp]
  \centering
  \includegraphics[width=0.72\textwidth]{phase3_eval_cm_t0_resnet50}
  \caption{Matrice de confusion normalisée --- ResNet50 (Tier~0, split de test).}
  \label{fig:phase3-cm-resnet50}
\end{figure}

Pour l'interprétabilité, le notebook génère des cartes Grad-CAM
(Figure~\ref{fig:phase3-gradcam}). ResNet50 (Tier~0) et CNN+HOG (Tier~1) focalisent
l'attention sur la région tumorale, validant visuellement la pertinence des représentations
apprises.

\begin{figure}[htbp]
  \centering
  \includegraphics[width=0.48\textwidth]{phase3_gradcam_t0_resnet50}\hfill
  \includegraphics[width=0.48\textwidth]{phase3_gradcam_t1_cnn_hog}
  \caption{Grad-CAM sur un cas correctement classifié par modèle~: ResNet50 (gauche) et
    CNN+HOG (droite).}
  \label{fig:phase3-gradcam}
\end{figure}

% -----------------------------------------------------------------------
\section{Exploration VLM (\texttt{VLM.ipynb})}

\subsection{Objectif et architecture}

Ce second notebook explore un \textbf{modèle vision-langage} appliqué à la même tâche de
classification en trois classes (Glioma, Meningioma, Pituitary Tumor). Le modèle retenu est
\textbf{Qwen2.5-VL-7B-Instruct}, chargé en \textbf{quantification 4-bit NF4} pour tenir sur une
GPU $\sim$16\,GiB, puis adapté par \textbf{LoRA} ($r=8$, $\alpha=16$, $\sim$0{,}3\,\% de
paramètres entraînables sur 8,3\,G de poids gelés).

L'image et une consigne textuelle sont fournies au modèle~; la réponse attendue est le nom de
la classe tumorale. Le jeu est découpé en 1\,050 / 150 / 300 échantillons (train / val / test).

\subsection{Teacher Forcing et sondes de probabilité}

Le notebook se concentre sur deux mécanismes fondamentaux du fine-tuning VLM~:

\begin{itemize}[leftmargin=1.4em,itemsep=2pt]
  \item \textbf{Teacher Forcing}~: pendant l'entraînement, les tokens de la réponse correcte
        sont fournis comme contexte~; un masque de labels restreint la perte aux seuls tokens
        de l'assistant (Figure~\ref{fig:vlm-tf-mask}).
  \item \textbf{Sondes de probabilité}~: avant et après LoRA, le notebook mesure la
        probabilité assignée à chaque token de classe à la première position de génération
        (Figures~\ref{fig:vlm-zero-shot} et~\ref{fig:vlm-before-after}).
\end{itemize}

Un \emph{QA Monitor} interroge le modèle à intervalles réguliers pendant l'entraînement
(5~époques, batch size~1, learning rate $2 \times 10^{-4}$) pour suivre l'évolution des
probabilités de classe (Figure~\ref{fig:vlm-training}).

\begin{figure}[htbp]
  \centering
  \includegraphics[width=0.88\textwidth]{teacher_forcing_mask}
  \caption{Masque de labels en mode Teacher Forcing~: seuls les tokens de la réponse assistant
    contribuent à la perte.}
  \label{fig:vlm-tf-mask}
\end{figure}

\begin{figure}[htbp]
  \centering
  \includegraphics[width=0.92\textwidth]{zero_shot_class_probs}
  \caption{Sonde zero-shot~: probabilités des tokens de classe avant fine-tuning LoRA
    (un échantillon par classe).}
  \label{fig:vlm-zero-shot}
\end{figure}

\begin{figure}[htbp]
  \centering
  \includegraphics[width=0.48\textwidth]{training_loss_curve}\hfill
  \includegraphics[width=0.48\textwidth]{qa_monitor_accuracy}
  \caption{Courbes d'entraînement (gauche) et précision du QA Monitor sur les échantillons-sondes
    (droite).}
  \label{fig:vlm-training}
\end{figure}

\begin{figure}[htbp]
  \centering
  \includegraphics[width=0.48\textwidth]{before_after_class_probs}\hfill
  \includegraphics[width=0.48\textwidth]{before_after_per_token_prob}
  \caption{Probabilités des tokens de classe avant vs après LoRA (gauche) et évolution
    per-token en teacher-forcing (droite).}
  \label{fig:vlm-before-after}
\end{figure}

\begin{figure}[htbp]
  \centering
  \includegraphics[width=0.88\textwidth]{before_after_loss_heatmap}
  \caption{Heatmap de la perte cross-entropie (teacher-forcing) avant vs après fine-tuning.}
  \label{fig:vlm-loss-heatmap}
\end{figure}

\subsection{Résultats}

Avant fine-tuning, le modèle généraliste distribue les probabilités de façon peu discriminante
sur les classes tumorales. Après LoRA, les sondes montrent un recentrage net sur le token de
la classe correcte, et le monitor atteint 3/3 sur les échantillons-sondes dès les premières
étapes d'évaluation.

Sur le jeu de test ($n = 300$), la génération greedy atteint une \textbf{accuracy globale de
94,7\,\%} (F1 par classe $\approx 0{,}95$). Le Tableau~\ref{tab:vlm-test} et la
Figure~\ref{fig:vlm-cm} détaillent ces scores.

\begin{table}[htbp]
  \centering
  \caption{Résultats VLM (Qwen2.5-VL-7B + LoRA) sur le jeu de test.}
  \label{tab:vlm-test}
  \begin{tabular}{@{}lccc@{}}
    \toprule
    \thd{Classe} & \thd{Precision} & \thd{Recall} & \thd{F1} \\
    \midrule
    Glioma          & 0,941 & 0,950 & 0,945 \\
    Meningioma      & 0,959 & 0,940 & 0,949 \\
    Pituitary Tumor & 0,960 & 0,950 & 0,955 \\
    \midrule
    \textbf{Accuracy globale} & & & \textbf{0,947} \\
    \bottomrule
  \end{tabular}
\end{table}

\begin{figure}[htbp]
  \centering
  \includegraphics[width=0.72\textwidth]{qwen2vl_lora_confusion_matrix}
  \caption{Matrice de confusion --- Qwen2.5-VL-7B + LoRA (jeu de test, $n = 300$).}
  \label{fig:vlm-cm}
\end{figure}

Ce résultat dépasse le benchmark ML (92,0\,\%) et le meilleur CNN pur (93,5\,\%), tout en
ouvrant la voie à des sorties textuelles explicatives --- avantage propre aux VLMs par rapport
aux classifieurs traditionnels. L'exploration reste toutefois préliminaire (un seul modèle,
fine-tuning léger, pas de validation statistique formelle).

% -----------------------------------------------------------------------
\section{Synthèse}

Les deux notebooks confirment que des représentations apprises (CNN, VLM) peuvent égaler ou
dépasser le pipeline manuel sur ce dataset, en particulier lorsque les branches hybrides
combinent CNN et descripteurs optimisés. Le VLM apporte en outre une capacité générative
(classification + justification textuelle) absente des approches ML et CNN classiques.

Les limites demeurent~: dépendance GPU, faible volume de données, absence de segmentation 3D
et protocole d'évaluation moins exhaustif que la Phase~2. Une fusion ML + DL + VLM, avec
validation multicentrique, constituerait la suite logique de ce travail.
